% GENERATED FILE. Edit canonical lessons, then run npm run generate:books.
% canonical-source-hash: fnv1a64:2e1f6bb7b031c815
% canonical-lessons: GE-C487-hausarzt, GE-C487-facharzt, GE-C487-untersuchen, GE-C487-behandeln, GE-C487-verschreiben

\chapter{Beim Hausarzt}
\label{ch:ge-beim-hausarzt}
\hlchaptermodality{487}

\begin{chapteropening}
\textbf{By the end of this chapter:} \emph{I can talk about seeing a doctor.}

Complete the last lesson of chapter 487: I can talk about seeing a doctor.
\end{chapteropening}

\section[der Hausarzt]{der Hausarzt — family doctor, GP}

\label{lesson:GE-C487-hausarzt}



\begin{quote}
Before the new one: say the German for to snore, then the German for nightmare.
\end{quote}

\subsection*{You'll want to know: der Hausarzt}
\textbf{der Hausarzt} — \textquotedblleft{}family doctor, GP\textquotedblright{}.

\begin{quote}
Mit Fieber gehe ich zuerst zum Hausarzt.
\end{quote}

\begin{quote}
\emph{With a fever I go to my GP first.}
\end{quote}

Plural: \emph{die Hausärzte}. A woman is \emph{die Hausärztin}.

\begin{cousinweb}
\emph{Hausarzt} is built on \emph{Haus} \textquotedblleft{}house\textquotedblright{} and \emph{Arzt} \textquotedblleft{}doctor\textquotedblright{}, which comes through Old High German \emph{arzāt} from Late Latin \emph{archiāter} \textquotedblleft{}physician\textquotedblright{} (EtymDB), a word from Greek.
\end{cousinweb}

\begin{grammarlens}[title={What you've built}]
One more word for telling, explaining and describing at B1.
\end{grammarlens}

\subsection*{Your turn}
\begin{itemize}
\raggedright
  \item \emph{Say it:} \emph{der Hausarzt}
  \item \emph{Say it:} \emph{der Hausarzt}, once more
  \item \emph{Say it:} a sentence of your own with \emph{der Hausarzt}
  \item \emph{Recall:} say the German for to snore, then the German for nightmare, then say \emph{der Hausarzt} again
\end{itemize}

\begin{tcolorbox}[breakable,colback=teal!4,colframe=teal!35!black,title={Before you move on}]
What does der Hausarzt mean? (\textquotedblleft{}Family doctor, GP\textquotedblright{}.) Say it once more.
\end{tcolorbox}
\section[der Facharzt]{der Facharzt — specialist (doctor)}

\label{lesson:GE-C487-facharzt}



\begin{quote}
Before the new one: say the German for nightmare, then the German for family doctor.
\end{quote}

\subsection*{You'll want to know: der Facharzt}
\textbf{der Facharzt} — \textquotedblleft{}specialist (doctor)\textquotedblright{}.

\begin{quote}
Der Hausarzt schickt mich zu einem Facharzt.
\end{quote}

\begin{quote}
\emph{The GP is sending me to a specialist.}
\end{quote}

Plural: \emph{die Fachärzte}. \emph{Das Fach} — subject, field.

\begin{cousinweb}
\emph{Facharzt} is built on \emph{Fach} \textquotedblleft{}compartment, field\textquotedblright{}, inherited from Old High German \emph{fah}, and \emph{Arzt} \textquotedblleft{}doctor\textquotedblright{}, from Late Latin \emph{archiāter} (EtymDB).
\end{cousinweb}

\begin{grammarlens}[title={What you've built}]
One more word for telling, explaining and describing at B1.
\end{grammarlens}

\subsection*{Your turn}
\begin{itemize}
\raggedright
  \item \emph{Say it:} \emph{der Facharzt}
  \item \emph{Say it:} \emph{der Facharzt}, once more
  \item \emph{Say it:} a sentence of your own with \emph{der Facharzt}
  \item \emph{Recall:} say the German for nightmare, then the German for family doctor, then say \emph{der Facharzt} again
\end{itemize}

\begin{tcolorbox}[breakable,colback=teal!4,colframe=teal!35!black,title={Before you move on}]
What does der Facharzt mean? (\textquotedblleft{}Specialist (doctor)\textquotedblright{}.) Say it once more.
\end{tcolorbox}
\section[untersuchen]{untersuchen — to examine}

\label{lesson:GE-C487-untersuchen}



\begin{quote}
Before the new one: say the German for family doctor, then the German for specialist.
\end{quote}

\subsection*{You'll want to know: untersuchen}
\textbf{untersuchen} — \textquotedblleft{}to examine\textquotedblright{}.

\begin{quote}
Die Ärztin untersucht mein Knie.
\end{quote}

\begin{quote}
\emph{The doctor examines my knee.}
\end{quote}

Inseparable: \emph{hat untersucht}. The noun: \emph{die Untersuchung}.

\begin{cousinweb}
\emph{untersuchen} is built on \emph{unter} \textquotedblleft{}under\textquotedblright{} and \emph{suchen} \textquotedblleft{}to seek\textquotedblright{}, inherited from Old High German \emph{suohhen}, from Proto-Germanic \emph{sōkijaną} (EtymDB).
\end{cousinweb}

\begin{grammarlens}[title={What you've built}]
One more word for telling, explaining and describing at B1.
\end{grammarlens}

\subsection*{Your turn}
\begin{itemize}
\raggedright
  \item \emph{Say it:} \emph{untersuchen}
  \item \emph{Say it:} \emph{untersuchen}, once more
  \item \emph{Say it:} a sentence of your own with \emph{untersuchen}
  \item \emph{Recall:} say the German for family doctor, then the German for specialist, then say \emph{untersuchen} again
\end{itemize}

\begin{tcolorbox}[breakable,colback=teal!4,colframe=teal!35!black,title={Before you move on}]
What does untersuchen mean? (\textquotedblleft{}To examine\textquotedblright{}.) Say it once more.
\end{tcolorbox}
\section[behandeln]{behandeln — to treat}

\label{lesson:GE-C487-behandeln}



\begin{quote}
Before the new one: say the German for specialist, then the German for to examine.
\end{quote}

\subsection*{You'll want to know: behandeln}
\textbf{behandeln} — \textquotedblleft{}to treat\textquotedblright{}.

\begin{quote}
Ein junger Arzt hat meinen Vater gut behandelt.
\end{quote}

\begin{quote}
\emph{A young doctor treated my father well.}
\end{quote}

Also of people in general: \emph{Er behandelt seine Kollegen fair.}

\begin{cousinweb}
\emph{behandeln} is built on \emph{be-} and \emph{handeln} \textquotedblleft{}to act, to handle\textquotedblright{}, from \emph{Hand}; EtymDB sets it beside Dutch \emph{behandelen} \textquotedblleft{}to handle, to treat\textquotedblright{}.
\end{cousinweb}

\begin{grammarlens}[title={What you've built}]
One more word for telling, explaining and describing at B1.
\end{grammarlens}

\subsection*{Your turn}
\begin{itemize}
\raggedright
  \item \emph{Say it:} \emph{behandeln}
  \item \emph{Say it:} \emph{behandeln}, once more
  \item \emph{Say it:} a sentence of your own with \emph{behandeln}
  \item \emph{Recall:} say the German for specialist, then the German for to examine, then say \emph{behandeln} again
\end{itemize}

\begin{tcolorbox}[breakable,colback=teal!4,colframe=teal!35!black,title={Before you move on}]
What does behandeln mean? (\textquotedblleft{}To treat\textquotedblright{}.) Say it once more.
\end{tcolorbox}
\section[verschreiben]{verschreiben — to prescribe}

\label{lesson:GE-C487-verschreiben}



\begin{quote}
Before the new one: say the German for to examine, then the German for to treat.
\end{quote}

\subsection*{You'll want to know: verschreiben}
\textbf{verschreiben} — \textquotedblleft{}to prescribe\textquotedblright{}.

\begin{quote}
Der Arzt hat mir Tabletten gegen den Husten verschrieben.
\end{quote}

\begin{quote}
\emph{The doctor prescribed me tablets for my cough.}
\end{quote}

Irregular: \emph{verschrieb, hat verschrieben}. \emph{Sich verschreiben} — to make a slip of the pen.

\begin{cousinweb}
\emph{verschreiben} is built on \emph{ver-} and \emph{schreiben}, which is inherited from Old High German \emph{scrīban} (EtymDB), an early borrowing from Latin \emph{scrībere} \textquotedblleft{}to write\textquotedblright{}.
\end{cousinweb}

\begin{grammarlens}[title={What you've built}]
One more word for telling, explaining and describing at B1.
\end{grammarlens}

\subsection*{Your turn}
\begin{itemize}
\raggedright
  \item \emph{Say it:} \emph{verschreiben}
  \item \emph{Say it:} \emph{verschreiben}, once more
  \item \emph{Say it:} a sentence of your own with \emph{verschreiben}
  \item \emph{Recall:} say the German for to examine, then the German for to treat, then say \emph{verschreiben} again
\end{itemize}

\begin{tcolorbox}[breakable,colback=teal!4,colframe=teal!35!black,title={Before you move on}]
What does verschreiben mean? (\textquotedblleft{}To prescribe\textquotedblright{}.) Say it once more.
\end{tcolorbox}
